% Options for packages loaded elsewhere
\PassOptionsToPackage{unicode}{hyperref}
\PassOptionsToPackage{hyphens}{url}
\documentclass[
  british,
  a4paper,
]{article}
\usepackage{xcolor}
\usepackage[margin=1in]{geometry}
\usepackage{amsmath,amssymb}
\setcounter{secnumdepth}{-\maxdimen} % remove section numbering
\usepackage{iftex}
\ifPDFTeX
  \usepackage[T1]{fontenc}
  \usepackage[utf8]{inputenc}
  \usepackage{textcomp} % provide euro and other symbols
\else % if luatex or xetex
  \usepackage{unicode-math} % this also loads fontspec
  \defaultfontfeatures{Scale=MatchLowercase}
  \defaultfontfeatures[\rmfamily]{Ligatures=TeX,Scale=1}
\fi
\usepackage{lmodern}
\ifPDFTeX\else
  % xetex/luatex font selection
  \setmainfont[]{Noto Serif}
  \setmonofont[]{Noto Sans Mono}
\fi
% Use upquote if available, for straight quotes in verbatim environments
\IfFileExists{upquote.sty}{\usepackage{upquote}}{}
\IfFileExists{microtype.sty}{% use microtype if available
  \usepackage[]{microtype}
  \UseMicrotypeSet[protrusion]{basicmath} % disable protrusion for tt fonts
}{}
\makeatletter
\@ifundefined{KOMAClassName}{% if non-KOMA class
  \IfFileExists{parskip.sty}{%
    \usepackage{parskip}
  }{% else
    \setlength{\parindent}{0pt}
    \setlength{\parskip}{6pt plus 2pt minus 1pt}}
}{% if KOMA class
  \KOMAoptions{parskip=half}}
\makeatother
% definitions for citeproc citations
\NewDocumentCommand\citeproctext{}{}
\NewDocumentCommand\citeproc{mm}{%
  \begingroup\def\citeproctext{#2}\cite{#1}\endgroup}
\makeatletter
 % allow citations to break across lines
 \let\@cite@ofmt\@firstofone
 % avoid brackets around text for \cite:
 \def\@biblabel#1{}
 \def\@cite#1#2{{#1\if@tempswa , #2\fi}}
\makeatother
\newlength{\cslhangindent}
\setlength{\cslhangindent}{1.5em}
\newlength{\csllabelwidth}
\setlength{\csllabelwidth}{3em}
\newenvironment{CSLReferences}[2] % #1 hanging-indent, #2 entry-spacing
 {\begin{list}{}{%
  \setlength{\itemindent}{0pt}
  \setlength{\leftmargin}{0pt}
  \setlength{\parsep}{0pt}
  % turn on hanging indent if param 1 is 1
  \ifodd #1
   \setlength{\leftmargin}{\cslhangindent}
   \setlength{\itemindent}{-1\cslhangindent}
  \fi
  % set entry spacing
  \setlength{\itemsep}{#2\baselineskip}}}
 {\end{list}}
\usepackage{calc}
\newcommand{\CSLBlock}[1]{\hfill\break\parbox[t]{\linewidth}{\strut\ignorespaces#1\strut}}
\newcommand{\CSLLeftMargin}[1]{\parbox[t]{\csllabelwidth}{\strut#1\strut}}
\newcommand{\CSLRightInline}[1]{\parbox[t]{\linewidth - \csllabelwidth}{\strut#1\strut}}
\newcommand{\CSLIndent}[1]{\hspace{\cslhangindent}#1}
\ifLuaTeX
\usepackage[bidi=basic,shorthands=off]{babel}
\else
\usepackage[bidi=default,shorthands=off]{babel}
\fi
\ifPDFTeX
\else
\babelfont{rm}[]{Noto Serif}
\fi
\ifLuaTeX
  \usepackage{selnolig} % disable illegal ligatures
\fi
\setlength{\emergencystretch}{3em} % prevent overfull lines
\providecommand{\tightlist}{%
  \setlength{\itemsep}{0pt}\setlength{\parskip}{0pt}}
\usepackage{array}
\usepackage{tabularx}
\usepackage{bookmark}
\IfFileExists{xurl.sty}{\usepackage{xurl}}{} % add URL line breaks if available
\urlstyle{same}
\hypersetup{
  pdflang={en-GB},
  hidelinks,
  pdfcreator={LaTeX via pandoc}}

\author{}
\date{}

\begin{document}

\section{Regular-entry prose pilot
03}\label{regular-entry-prose-pilot-03}

\emph{Assembly-only pilot for revised compact regular-entry commentary
in a more direct factual style. Trace rendering is reused, and the
commentary is supplied from a separate book-prose layer.}

\subsection{bake --- OE bacan}\label{bake-oe-bacan}

Derivation: \emph{*bákaną} \textgreater{} \emph{bacan} (regular).

\begingroup
\setlength{\fboxsep}{6pt}

\noindent\fbox{%
\begin{minipage}{0.97\linewidth}
\small
\begin{tabularx}{\linewidth}{@{}>{\raggedright\arraybackslash}p{0.300\linewidth}>{\centering\arraybackslash}X>{\raggedright\arraybackslash}p{0.540\linewidth}@{}}
\begin{minipage}[t]{\linewidth}
\centering\textbf{Earlier Germanic changes}\par
\vspace{0.35em}
\raggedright
\centering\textbf{West Germanic}\par
\raggedright
\vspace{0.2em}
\raggedright [no change]\par
\vspace{0.6em}
\centering\textbf{Northwest Germanic}\par
\raggedright
\vspace{0.2em}
\raggedright [no change]\par
\end{minipage}
&
&
\begin{minipage}[t]{\linewidth}
\centering\textbf{Old English changes}\par
\vspace{0.35em}
\raggedright
\begin{tabular}{@{}>{\raggedright\arraybackslash}p{0.74\linewidth}@{\hspace{0.55em}}>{\raggedright\arraybackslash}p{0.16\linewidth}@{\hspace{0.25em}}}
Anglo Frisian Brightening & \emph{*bækaną} \\
OE A Restoration & \emph{*bakaną} \\
OE Heavy Syllable Nasal Apocope & \emph{*bakan} \\
OE Secondary Nasalization & \emph{*bakąn} \\
OE Weak Tail Reduction & \emph{*bakan} \\
\end{tabular}
\end{minipage}
\\
\end{tabularx}
\end{minipage}%
} \endgroup

Orel reconstructs \emph{*bakanan}, and Campbell cites \emph{bacan} as a
standard example of A-restoration (\citeproc{ref-Orel2003}{Orel 2003};
\citeproc{ref-Campbell1959}{Campbell 1959, 61};
\citeproc{ref-RingeTaylor2014}{Ringe and Taylor 2014, 270}).
Bosworth-Toller and Clark Hall record \emph{bacan} as the Old English
infinitive (\citeproc{ref-BosworthToller1898}{Bosworth and Toller 1898,
72}; \citeproc{ref-ClarkHall1960}{Clark Hall 1960, 19}). The selected
input \emph{*bákaną} represents the same weak verb in the accent
notation used here. From \emph{*bákaną}, Anglo-Frisian brightening first
gives \emph{*bækaną}. A-restoration before single \emph{k} returns the
stem vowel to \emph{a}, and later weak-tail reductions yield
\emph{bacan}.

\subsection{beech --- OE bōc}\label{beech-oe-bux14dc}

Derivation: \emph{*bōkō} \textgreater{} \emph{bōc} (regular).

\begingroup
\setlength{\fboxsep}{6pt}

\noindent\fbox{%
\begin{minipage}{0.97\linewidth}
\small
\begin{tabularx}{\linewidth}{@{}>{\raggedright\arraybackslash}p{0.460\linewidth}>{\centering\arraybackslash}X>{\raggedright\arraybackslash}p{0.460\linewidth}@{}}
\begin{minipage}[t]{\linewidth}
\centering\textbf{Earlier Germanic changes}\par
\vspace{0.35em}
\raggedright
\centering\textbf{West Germanic}\par
\raggedright
\vspace{0.2em}
\raggedright [no change]\par
\vspace{0.6em}
\centering\textbf{Northwest Germanic}\par
\raggedright
\vspace{0.2em}
\begin{tabular}{@{}>{\raggedright\arraybackslash}p{0.74\linewidth}@{\hspace{0.55em}}>{\raggedright\arraybackslash}p{0.16\linewidth}@{\hspace{0.25em}}}
\mbox{NWGmc Final Long O Raising} & \emph{*bōku} \\
\end{tabular}
\end{minipage}
&
&
\begin{minipage}[t]{\linewidth}
\centering\textbf{Old English changes}\par
\vspace{0.35em}
\raggedright
\begin{tabular}{@{}>{\raggedright\arraybackslash}p{0.74\linewidth}@{\hspace{0.55em}}>{\raggedright\arraybackslash}p{0.16\linewidth}@{\hspace{0.25em}}}
\mbox{OE High Vowel Apocope} & \emph{*bōk} \\
\end{tabular}
\end{minipage}
\\
\end{tabularx}
\end{minipage}%
} \endgroup

Kroonen gives \emph{*bōk(j)ō-} with Old English \emph{boc} and oblique
\emph{bēce} (\citeproc{ref-Kroonen2013}{Kroonen 2013, 81}). The selected
input \emph{*bōkō} is the nominative-singular shape of that family. The
selected form here is nominative \emph{bōc}; \emph{bēċe} belongs to the
same paradigm but is not the citation form used for this derivation.
From \emph{*bōkō}, final long \emph{ō} raises in Northwest Germanic, and
high-vowel apocope yields \emph{bōc}.

\subsection{bier --- OE bǣr}\label{bier-oe-bux1e3r}

Derivation: \emph{*bḗrō} \textgreater{} \emph{bǣr} (regular).

\begingroup
\setlength{\fboxsep}{6pt}

\noindent\fbox{%
\begin{minipage}{0.97\linewidth}
\small
\begin{tabularx}{\linewidth}{@{}>{\raggedright\arraybackslash}p{0.460\linewidth}>{\centering\arraybackslash}X>{\raggedright\arraybackslash}p{0.460\linewidth}@{}}
\begin{minipage}[t]{\linewidth}
\centering\textbf{Earlier Germanic changes}\par
\vspace{0.35em}
\raggedright
\centering\textbf{West Germanic}\par
\raggedright
\vspace{0.2em}
\raggedright [no change]\par
\vspace{0.6em}
\centering\textbf{Northwest Germanic}\par
\raggedright
\vspace{0.2em}
\begin{tabular}{@{}>{\raggedright\arraybackslash}p{0.74\linewidth}@{\hspace{0.55em}}>{\raggedright\arraybackslash}p{0.16\linewidth}@{\hspace{0.25em}}}
\mbox{NWGmc Final Long O Raising} & \emph{*bḗru} \\
\mbox{NWGmc Long E Lowering} & \emph{*bǣru} \\
\end{tabular}
\end{minipage}
&
&
\begin{minipage}[t]{\linewidth}
\centering\textbf{Old English changes}\par
\vspace{0.35em}
\raggedright
\begin{tabular}{@{}>{\raggedright\arraybackslash}p{0.74\linewidth}@{\hspace{0.55em}}>{\raggedright\arraybackslash}p{0.16\linewidth}@{\hspace{0.25em}}}
\mbox{OE High Vowel Apocope} & \emph{*bǣr} \\
\end{tabular}
\end{minipage}
\\
\end{tabularx}
\end{minipage}%
} \endgroup

Kroonen reconstructs \emph{*bērō-} and cites Old English \emph{bar},
\emph{bær} among the reflexes (\citeproc{ref-Kroonen2013}{Kroonen 2013,
717}; \citeproc{ref-ClarkHall1960}{Clark Hall 1960, 32};
\citeproc{ref-BosworthToller1898}{Bosworth and Toller 1898, 73}). The
selected input \emph{*bḗrō} is the same noun in the accent notation used
here. Clark Hall and Bosworth-Toller give \emph{bær}, and Kroonen also
records \emph{bar}. The normalized form \emph{bǣr} marks the long vowel
explicitly. From \emph{*bḗrō}, final long \emph{ō} raises, long \emph{ē}
lowers to \emph{ǣ}, and apocope yields \emph{bǣr}.

\subsection{both --- OE bū}\label{both-oe-bux16b}

Derivation: \emph{*bō} \textgreater{} \emph{bū} (regular).

\begingroup
\setlength{\fboxsep}{6pt}

\noindent\fbox{%
\begin{minipage}{0.97\linewidth}
\small
\begin{tabularx}{\linewidth}{@{}>{\raggedright\arraybackslash}p{0.540\linewidth}>{\centering\arraybackslash}X>{\raggedright\arraybackslash}p{0.300\linewidth}@{}}
\begin{minipage}[t]{\linewidth}
\centering\textbf{Earlier Germanic changes}\par
\vspace{0.35em}
\raggedright
\centering\textbf{West Germanic}\par
\raggedright
\vspace{0.2em}
\raggedright [no change]\par
\vspace{0.6em}
\centering\textbf{Northwest Germanic}\par
\raggedright
\vspace{0.2em}
\begin{tabular}{@{}>{\raggedright\arraybackslash}p{0.74\linewidth}@{\hspace{0.55em}}>{\raggedright\arraybackslash}p{0.16\linewidth}@{\hspace{0.25em}}}
NWGmc Stressed Monosyllable O Raising & \emph{*bū} \\
\end{tabular}
\end{minipage}
&
&
\begin{minipage}[t]{\linewidth}
\centering\textbf{Old English changes}\par
\vspace{0.35em}
\raggedright
\raggedright [no change]\par
\end{minipage}
\\
\end{tabularx}
\end{minipage}%
} \endgroup

Kroonen treats the numeral under \emph{*ba-} and includes the inherited
\emph{*bō} beside the extended dual forms
(\citeproc{ref-Kroonen2013}{Kroonen 2013, 47}). Brunner, Campbell, and
Fulk preserve the Old English dual paradigm with masculine \emph{bēġen},
feminine \emph{bā}, and neuter \emph{bū}
(\citeproc{ref-SieversBrunner1965}{Sievers 1965, 324};
\citeproc{ref-Campbell1959}{Campbell 1959, 295};
\citeproc{ref-Fulk2018}{Fulk 2018, 223}). The selected form is neuter
\emph{bū}. Campbell and Brunner treat stressed monosyllabic \emph{ō
\textgreater{} ū} as a regular development in forms of this type, and
\emph{bū} belongs to that set (\citeproc{ref-Campbell1959}{Campbell
1959, 122}; \citeproc{ref-SieversBrunner1965}{Sievers 1965, 69}).

Comparison note. Orel cites an older \emph{*bō-j-} explanation for
\emph{bēġen}, while Fulk reports Seebold's preference for \emph{*bō-þ-}
(\citeproc{ref-Orel2003}{Orel 2003, 45}; \citeproc{ref-Fulk2018}{Fulk
2018, 223}). Those analyses concern \emph{bēġen} and the related
extended forms. The derivation here is the neuter \emph{*bō
\textgreater{} bū}.

\subsection{bow --- OE bīeġan}\label{bow-oe-bux12beux121an}

Derivation: \emph{*báugijaną} \textgreater{} \emph{bīeġan} (regular).

\begingroup
\setlength{\fboxsep}{6pt}

\noindent\fbox{%
\begin{minipage}{0.97\linewidth}
\footnotesize
\begin{tabularx}{\linewidth}{@{}>{\raggedright\arraybackslash}p{0.280\linewidth}>{\centering\arraybackslash}X>{\raggedright\arraybackslash}p{0.560\linewidth}@{}}
\begin{minipage}[t]{\linewidth}
\centering\textbf{Earlier Germanic changes}\par
\vspace{0.35em}
\raggedright
\centering\textbf{West Germanic}\par
\raggedright
\vspace{0.2em}
\raggedright [no change]\par
\vspace{0.6em}
\centering\textbf{Northwest Germanic}\par
\raggedright
\vspace{0.2em}
\raggedright [no change]\par
\end{minipage}
&
&
\begin{minipage}[t]{\linewidth}
\centering\textbf{Old English changes}\par
\vspace{0.35em}
\raggedright
\begin{tabular}{@{}>{\raggedright\arraybackslash}p{0.62\linewidth}@{\hspace{0.55em}}>{\raggedright\arraybackslash}p{0.28\linewidth}@{\hspace{0.25em}}}
OE Au Fronting & \emph{*báeugijaną} \\
OE Diphthong Leveling & \emph{*bēagijaną} \\
OE Heavy Syllable Nasal Apocope & \emph{*bēagijan} \\
OE Secondary Nasalization & \emph{*bēagijąn} \\
Sievers Law Syncope & \emph{*bēagjąn} \\
OE Velar Palatalization & \emph{*bēaʤjąn} \\
OE I Umlaut & \emph{*bīeʤjąn} \\
OE Weak Tail Reduction & \emph{*bīeʤjan} \\
OE J Loss After Heavy & \emph{*bīeʤan} \\
\end{tabular}
\end{minipage}
\\
\end{tabularx}
\end{minipage}%
} \endgroup

Kroonen reconstructs the weak causative \emph{*baugjan-}, and Ringe and
Taylor give the fuller development to West Saxon \emph{biegan}
(\citeproc{ref-Kroonen2013}{Kroonen 2013, 75};
\citeproc{ref-RingeTaylor2014}{Ringe and Taylor 2014, 94}). The verb
belongs to the weak causative member of the bend-family. Clark Hall and
Bosworth-Toller record the Old English weak verb
(\citeproc{ref-ClarkHall1960}{Clark Hall 1960, 34};
\citeproc{ref-BosworthToller1898}{Bosworth and Toller 1898, 102}). The
spelling \emph{bīeġan} used here normalizes that attested form. From
pre-Old-English \emph{*bēagjan}, palatalization of \emph{gj} and
i-mutation yield the West Saxon infinitive.

\subsection{fare --- OE faran}\label{fare-oe-faran}

Derivation: \emph{*fáraną} \textgreater{} \emph{faran} (regular).

\begingroup
\setlength{\fboxsep}{6pt}

\noindent\fbox{%
\begin{minipage}{0.97\linewidth}
\small
\begin{tabularx}{\linewidth}{@{}>{\raggedright\arraybackslash}p{0.300\linewidth}>{\centering\arraybackslash}X>{\raggedright\arraybackslash}p{0.540\linewidth}@{}}
\begin{minipage}[t]{\linewidth}
\centering\textbf{Earlier Germanic changes}\par
\vspace{0.35em}
\raggedright
\centering\textbf{West Germanic}\par
\raggedright
\vspace{0.2em}
\raggedright [no change]\par
\vspace{0.6em}
\centering\textbf{Northwest Germanic}\par
\raggedright
\vspace{0.2em}
\raggedright [no change]\par
\end{minipage}
&
&
\begin{minipage}[t]{\linewidth}
\centering\textbf{Old English changes}\par
\vspace{0.35em}
\raggedright
\begin{tabular}{@{}>{\raggedright\arraybackslash}p{0.74\linewidth}@{\hspace{0.55em}}>{\raggedright\arraybackslash}p{0.16\linewidth}@{\hspace{0.25em}}}
Anglo Frisian Brightening & \emph{*færaną} \\
OE A Restoration & \emph{*faraną} \\
OE Heavy Syllable Nasal Apocope & \emph{*faran} \\
OE Secondary Nasalization & \emph{*farąn} \\
OE Weak Tail Reduction & \emph{*faran} \\
\end{tabular}
\end{minipage}
\\
\end{tabularx}
\end{minipage}%
} \endgroup

Kroonen gives the inherited strong verb as \emph{*faran-}, and Orel
gives the same lexeme as \emph{*faranan}, both with Old English
\emph{faran} among the reflexes (\citeproc{ref-Kroonen2013}{Kroonen
2013, 149}; \citeproc{ref-Orel2003}{Orel 2003, 132};
\citeproc{ref-Campbell1959}{Campbell 1959, 61}). Clark Hall and
Bosworth-Toller distinguish strong \emph{faran} from weak \emph{færan},
while forms such as \emph{fære} and \emph{færð} belong to the present
tense (\citeproc{ref-ClarkHall1960}{Clark Hall 1960, 113};
\citeproc{ref-BosworthToller1898}{Bosworth and Toller 1898, 250}). From
\emph{*fáraną}, brightening first gives \emph{*færaną}; A-restoration
before single \emph{r} returns \emph{*faraną}, and later reductions
yield \emph{faran}.

\subsection{guest --- OE ġiest}\label{guest-oe-ux121iest}

Derivation: \emph{*gástiz} \textgreater{} \emph{ġiest} (regular).

\begingroup
\setlength{\fboxsep}{6pt}

\noindent\fbox{%
\begin{minipage}{0.97\linewidth}
\small
\begin{tabularx}{\linewidth}{@{}>{\raggedright\arraybackslash}p{0.460\linewidth}>{\centering\arraybackslash}X>{\raggedright\arraybackslash}p{0.460\linewidth}@{}}
\begin{minipage}[t]{\linewidth}
\centering\textbf{Earlier Germanic changes}\par
\vspace{0.35em}
\raggedright
\centering\textbf{West Germanic}\par
\raggedright
\vspace{0.2em}
\raggedright [no change]\par
\vspace{0.6em}
\centering\textbf{Northwest Germanic}\par
\raggedright
\vspace{0.2em}
\begin{tabular}{@{}>{\raggedright\arraybackslash}p{0.74\linewidth}@{\hspace{0.55em}}>{\raggedright\arraybackslash}p{0.16\linewidth}@{\hspace{0.25em}}}
\mbox{PGmc Final Z Deletion} & \emph{*gásti} \\
\end{tabular}
\end{minipage}
&
&
\begin{minipage}[t]{\linewidth}
\centering\textbf{Old English changes}\par
\vspace{0.35em}
\raggedright
\begin{tabular}{@{}>{\raggedright\arraybackslash}p{0.74\linewidth}@{\hspace{0.55em}}>{\raggedright\arraybackslash}p{0.16\linewidth}@{\hspace{0.25em}}}
Anglo Frisian Brightening & \emph{*gæsti} \\
OE Velar Palatalization & \emph{*ʤæsti} \\
OE I Umlaut & \emph{*ʤesti} \\
OE Ws Palatal Diphthongization & \emph{*ʤiesti} \\
\mbox{OE High Vowel Apocope} & \emph{*ʤiest} \\
\end{tabular}
\end{minipage}
\\
\end{tabularx}
\end{minipage}%
} \endgroup

Campbell and Ringe and Taylor treat this as an ordinary i-stem whose
West Saxon development yields \emph{ġiest}, while non-West-Saxon
evidence preserves forms of the \emph{gest} type
(\citeproc{ref-Campbell1959}{Campbell 1959, 91};
\citeproc{ref-RingeTaylor2014}{Ringe and Taylor 2014, 296}).
Bosworth-Toller and Clark Hall record spellings such as \emph{gist},
\emph{gest}, \emph{giest}, and \emph{gyst}
(\citeproc{ref-BosworthToller1898}{Bosworth and Toller 1898, 387};
\citeproc{ref-ClarkHall1960}{Clark Hall 1960, 149}). The derivation is
regular: brightening and i-mutation reshape the stem, and West Saxon
palatal diphthongization gives \emph{ie}.

Dialect note. Anglian \emph{gest} is genuine Old English evidence
alongside West Saxon \emph{ġiest} (\citeproc{ref-RingeTaylor2014}{Ringe
and Taylor 2014, 296}; \citeproc{ref-BosworthToller1898}{Bosworth and
Toller 1898, 387}).

\subsection{learn --- OE liornian}\label{learn-oe-liornian}

Derivation: \emph{*líznōjaną} \textgreater{} \emph{liornian} (regular).

\begingroup
\setlength{\fboxsep}{6pt}

\noindent\fbox{%
\begin{minipage}{0.97\linewidth}
\footnotesize
\begin{tabularx}{\linewidth}{@{}>{\raggedright\arraybackslash}p{0.280\linewidth}>{\centering\arraybackslash}X>{\raggedright\arraybackslash}p{0.560\linewidth}@{}}
\begin{minipage}[t]{\linewidth}
\centering\textbf{Earlier Germanic changes}\par
\vspace{0.35em}
\raggedright
\centering\textbf{West Germanic}\par
\raggedright
\vspace{0.2em}
\raggedright [no change]\par
\vspace{0.6em}
\centering\textbf{Northwest Germanic}\par
\raggedright
\vspace{0.2em}
\raggedright [no change]\par
\end{minipage}
&
&
\begin{minipage}[t]{\linewidth}
\centering\textbf{Old English changes}\par
\vspace{0.35em}
\raggedright
\begin{tabular}{@{}>{\raggedright\arraybackslash}p{0.62\linewidth}@{\hspace{0.55em}}>{\raggedright\arraybackslash}p{0.28\linewidth}@{\hspace{0.25em}}}
OE Breaking & \emph{*líornōjaną} \\
OE Heavy Syllable Nasal Apocope & \emph{*líornōjan} \\
OE Secondary Nasalization & \emph{*líornōjąn} \\
OE I Umlaut & \emph{*líornējąn} \\
OE Unstressed Long Vowel Shortening & \emph{*líornejąn} \\
OE Weak Tail Reduction & \emph{*líornejan} \\
OE Intervocalic J Vocalization & \emph{*líorneian} \\
OE Unstressed EI Contraction & \emph{*líornian} \\
\end{tabular}
\end{minipage}
\\
\end{tabularx}
\end{minipage}%
} \endgroup

Ringe and Taylor place the verb in a class-II weak family of the
\emph{*liznō-} type, and Kroonen keeps the same comparative base
(\citeproc{ref-RingeTaylor2014}{Ringe and Taylor 2014, 318};
\citeproc{ref-Kroonen2013}{Kroonen 2013, 328}). Campbell records
Northumbrian \emph{liornian} and treats the \emph{eo} of \emph{leornian}
as secondary (\citeproc{ref-Campbell1959}{Campbell 1959, 123};
\citeproc{ref-Campbell1959}{1959, 296}). The Northumbrian spelling
preserves \emph{io} where original \emph{eo} and \emph{io} remain
distinct. The selected target \emph{liornian} therefore represents the
Northumbrian regular form.

Dialect note. Clark Hall and Bright give \emph{leornian} as the usual
headword (\citeproc{ref-ClarkHall1960}{Clark Hall 1960, 193};
\citeproc{ref-BrightCassidyRingler1971}{Bright 1971, 58}). From
\emph{*líznōjaną}, the regular development yields \emph{liornian};
\emph{leornian} reflects the later \emph{eo} development.

\subsection{linden --- OE lind}\label{linden-oe-lind}

Derivation: \emph{*líndō} \textgreater{} \emph{lind} (regular).

\begingroup
\setlength{\fboxsep}{6pt}

\noindent\fbox{%
\begin{minipage}{0.97\linewidth}
\small
\begin{tabularx}{\linewidth}{@{}>{\raggedright\arraybackslash}p{0.460\linewidth}>{\centering\arraybackslash}X>{\raggedright\arraybackslash}p{0.460\linewidth}@{}}
\begin{minipage}[t]{\linewidth}
\centering\textbf{Earlier Germanic changes}\par
\vspace{0.35em}
\raggedright
\centering\textbf{West Germanic}\par
\raggedright
\vspace{0.2em}
\raggedright [no change]\par
\vspace{0.6em}
\centering\textbf{Northwest Germanic}\par
\raggedright
\vspace{0.2em}
\begin{tabular}{@{}>{\raggedright\arraybackslash}p{0.74\linewidth}@{\hspace{0.55em}}>{\raggedright\arraybackslash}p{0.16\linewidth}@{\hspace{0.25em}}}
\mbox{NWGmc Final Long O Raising} & \emph{*líndu} \\
\end{tabular}
\end{minipage}
&
&
\begin{minipage}[t]{\linewidth}
\centering\textbf{Old English changes}\par
\vspace{0.35em}
\raggedright
\begin{tabular}{@{}>{\raggedright\arraybackslash}p{0.74\linewidth}@{\hspace{0.55em}}>{\raggedright\arraybackslash}p{0.16\linewidth}@{\hspace{0.25em}}}
\mbox{OE High Vowel Apocope} & \emph{*línd} \\
\end{tabular}
\end{minipage}
\\
\end{tabularx}
\end{minipage}%
} \endgroup

Kroonen reconstructs \emph{*lindō-}, and Clark Hall and Bosworth-Toller
record Old English \emph{lind} `linden, lime-tree'
(\citeproc{ref-Kroonen2013}{Kroonen 2013, 337};
\citeproc{ref-ClarkHall1960}{Clark Hall 1960, 198};
\citeproc{ref-BosworthToller1898}{Bosworth and Toller 1898, 630}). The
selected input \emph{*líndō} is the citation-form noun used here.
Northwest Germanic final \emph{ō} raises to \emph{*líndu}, and
high-vowel apocope yields \emph{lind}. Clark Hall also gives adjectival
\emph{linden} `made of linden-wood'; that is a related adjective, not
the noun treated here (\citeproc{ref-ClarkHall1960}{Clark Hall 1960,
198}).

\subsection{mother --- OE mōder}\label{mother-oe-mux14dder}

Derivation: \emph{*mōdēr} \textgreater{} \emph{mōder} (regular).

\begingroup
\setlength{\fboxsep}{6pt}

\noindent\fbox{%
\begin{minipage}{0.97\linewidth}
\small
\begin{tabularx}{\linewidth}{@{}>{\raggedright\arraybackslash}p{0.460\linewidth}>{\centering\arraybackslash}X>{\raggedright\arraybackslash}p{0.460\linewidth}@{}}
\begin{minipage}[t]{\linewidth}
\centering\textbf{Earlier Germanic changes}\par
\vspace{0.35em}
\raggedright
\centering\textbf{West Germanic}\par
\raggedright
\vspace{0.2em}
\raggedright [no change]\par
\vspace{0.6em}
\centering\textbf{Northwest Germanic}\par
\raggedright
\vspace{0.2em}
\begin{tabular}{@{}>{\raggedright\arraybackslash}p{0.74\linewidth}@{\hspace{0.55em}}>{\raggedright\arraybackslash}p{0.16\linewidth}@{\hspace{0.25em}}}
\mbox{NWGmc Long E Lowering} & \emph{*mōdǣr} \\
\end{tabular}
\end{minipage}
&
&
\begin{minipage}[t]{\linewidth}
\centering\textbf{Old English changes}\par
\vspace{0.35em}
\raggedright
\begin{tabular}{@{}>{\raggedright\arraybackslash}p{0.74\linewidth}@{\hspace{0.55em}}>{\raggedright\arraybackslash}p{0.16\linewidth}@{\hspace{0.25em}}}
OE Unstressed Long Vowel Shortening & \emph{*mōdær} \\
OE Unstressed AE Merger & \emph{*mōder} \\
\end{tabular}
\end{minipage}
\\
\end{tabularx}
\end{minipage}%
} \endgroup

Kroonen and Orel cite the Proto-Germanic r-stem as \emph{*mōdēr-} or
\emph{*mōdēr} (\citeproc{ref-Kroonen2013}{Kroonen 2013, 368};
\citeproc{ref-Orel2003}{Orel 2003, 274}). Clark Hall, Campbell, and
Ringe and Taylor preserve the Old English headword tradition \emph{mōdor
/ modor} with oblique \emph{mēder} in the paradigm
(\citeproc{ref-ClarkHall1960}{Clark Hall 1960, 219};
\citeproc{ref-Campbell1959}{Campbell 1959, 598};
\citeproc{ref-RingeTaylor2014}{Ringe and Taylor 2014, 312}). The
selected form here is nominative \emph{mōder}. From \emph{*mōdēr},
regular suffixal development yields that form.

Comparison note. \emph{Mōdor} is the usual dictionary form and reflects
later levelling within the r-stem paradigm. \emph{Mēder} preserves the
oblique line.

\subsection{show --- OE sċēawian}\label{show-oe-sux10bux113awian}

Derivation: \emph{*skáwōjaną} \textgreater{} \emph{sċēawian} (regular).

\begingroup
\setlength{\fboxsep}{6pt}

\noindent\fbox{%
\begin{minipage}{0.97\linewidth}
\footnotesize
\begin{tabularx}{\linewidth}{@{}>{\raggedright\arraybackslash}p{0.280\linewidth}>{\centering\arraybackslash}X>{\raggedright\arraybackslash}p{0.560\linewidth}@{}}
\begin{minipage}[t]{\linewidth}
\centering\textbf{Earlier Germanic changes}\par
\vspace{0.35em}
\raggedright
\centering\textbf{West Germanic}\par
\raggedright
\vspace{0.2em}
\raggedright [no change]\par
\vspace{0.6em}
\centering\textbf{Northwest Germanic}\par
\raggedright
\vspace{0.2em}
\raggedright [no change]\par
\end{minipage}
&
&
\begin{minipage}[t]{\linewidth}
\centering\textbf{Old English changes}\par
\vspace{0.35em}
\raggedright
\begin{tabular}{@{}>{\raggedright\arraybackslash}p{0.62\linewidth}@{\hspace{0.55em}}>{\raggedright\arraybackslash}p{0.28\linewidth}@{\hspace{0.25em}}}
OE Aw Long Diphthong & \emph{*skḗawōjaną} \\
OE Heavy Syllable Nasal Apocope & \emph{*skḗawōjan} \\
OE Secondary Nasalization & \emph{*skḗawōjąn} \\
OE Sk Palatalization & \emph{*ʃḗawōjąn} \\
OE I Umlaut & \emph{*ʃḗawējąn} \\
OE Unstressed Long Vowel Shortening & \emph{*ʃḗawejąn} \\
OE Weak Tail Reduction & \emph{*ʃḗawejan} \\
OE Intervocalic J Vocalization & \emph{*ʃḗaweian} \\
OE Unstressed EI Contraction & \emph{*ʃḗawian} \\
\end{tabular}
\end{minipage}
\\
\end{tabularx}
\end{minipage}%
} \endgroup

Orel and Kroonen place the verb in the ordinary class-II show-family,
and Brunner records Old English \emph{scēawian}, \emph{scāwian}
(\citeproc{ref-Orel2003}{Orel 2003, 339};
\citeproc{ref-Kroonen2013}{Kroonen 2013, 482};
\citeproc{ref-SieversBrunner1965}{Sievers 1965, 148}). From
\emph{*skáwōjaną}, \emph{aw} before a following vowel yields Old English
\emph{ēaw}, while the class-II suffix preserves \emph{ō} between
\emph{w} and \emph{j}. The result is \emph{sċēawian}
(\citeproc{ref-Campbell1959}{Campbell 1959, 138};
\citeproc{ref-Orel2003}{Orel 2003, 339}).

Form note. The spelling \emph{sċēawian} normalizes initial in dictionary
\emph{scēawian} (\citeproc{ref-Campbell1959}{Campbell 1959, 19};
\citeproc{ref-Hogg1992}{Hogg 1992, 236}).

\subsection{weapon --- OE wǣpn}\label{weapon-oe-wux1e3pn}

Derivation: \emph{*wḗpną} \textgreater{} \emph{wǣpn} (regular).

\begingroup
\setlength{\fboxsep}{6pt}

\noindent\fbox{%
\begin{minipage}{0.97\linewidth}
\small
\begin{tabularx}{\linewidth}{@{}>{\raggedright\arraybackslash}p{0.460\linewidth}>{\centering\arraybackslash}X>{\raggedright\arraybackslash}p{0.460\linewidth}@{}}
\begin{minipage}[t]{\linewidth}
\centering\textbf{Earlier Germanic changes}\par
\vspace{0.35em}
\raggedright
\centering\textbf{West Germanic}\par
\raggedright
\vspace{0.2em}
\raggedright [no change]\par
\vspace{0.6em}
\centering\textbf{Northwest Germanic}\par
\raggedright
\vspace{0.2em}
\begin{tabular}{@{}>{\raggedright\arraybackslash}p{0.74\linewidth}@{\hspace{0.55em}}>{\raggedright\arraybackslash}p{0.16\linewidth}@{\hspace{0.25em}}}
\mbox{NWGmc Long E Lowering} & \emph{*wǣpną} \\
\end{tabular}
\end{minipage}
&
&
\begin{minipage}[t]{\linewidth}
\centering\textbf{Old English changes}\par
\vspace{0.35em}
\raggedright
\begin{tabular}{@{}>{\raggedright\arraybackslash}p{0.74\linewidth}@{\hspace{0.55em}}>{\raggedright\arraybackslash}p{0.16\linewidth}@{\hspace{0.25em}}}
OE Heavy Syllable Nasal Apocope & \emph{*wǣpn} \\
\end{tabular}
\end{minipage}
\\
\end{tabularx}
\end{minipage}%
} \endgroup

Kroonen reconstructs a double-stem noun \emph{*wēbna- \textasciitilde{}
wēpna-} and cites Old English \emph{wæpn} among the reflexes
(\citeproc{ref-Kroonen2013}{Kroonen 2013, 617};
\citeproc{ref-Campbell1959}{Campbell 1959, 150};
\citeproc{ref-Campbell1959}{1959, 226}). Campbell preserves unbroken
\emph{wépn} beside broken \emph{wépen}-type forms. Bright contrasts
broken \emph{wǣpen/wapen} with unbroken \emph{wǣpnes}, and Clark Hall
lemmatizes the noun under \emph{wapen} while preserving unbroken forms
elsewhere (\citeproc{ref-BrightCassidyRingler1971}{Bright 1971, 29};
\citeproc{ref-ClarkHall1960}{Clark Hall 1960, 355}). The selected form
\emph{wǣpn} represents the unbroken line from \emph{*wḗpną}.

Form note. \emph{Wǣpen} is the usual late West Saxon simplex headword,
while \emph{wǣpnes} and related forms preserve the unbroken cluster
(\citeproc{ref-Campbell1959}{Campbell 1959, 150};
\citeproc{ref-Campbell1959}{1959, 226};
\citeproc{ref-BrightCassidyRingler1971}{Bright 1971, 29};
\citeproc{ref-ClarkHall1960}{Clark Hall 1960, 355}).

\clearpage

\subsection*{References}\label{references}
\addcontentsline{toc}{subsection}{References}

\protect\phantomsection\label{refs}
\begin{CSLReferences}{1}{1}
\bibitem[\citeproctext]{ref-BosworthToller1898}
Bosworth, Joseph, and T. Northcote Toller. 1898. \emph{An {Anglo-Saxon}
Dictionary}. Clarendon Press.

\bibitem[\citeproctext]{ref-BrightCassidyRingler1971}
Bright, James W. 1971. \emph{Bright's {Old English} Grammar and Reader}.
3rd edn. Edited by Frederic G. Cassidy and Richard N. Ringler. Holt,
Rinehart; Winston.

\bibitem[\citeproctext]{ref-Campbell1959}
Campbell, Alistair. 1959. \emph{Old {English} Grammar}. Clarendon Press.

\bibitem[\citeproctext]{ref-ClarkHall1960}
Clark Hall, J. R. 1960. \emph{A Concise {Anglo-Saxon} Dictionary}. 4th
edn. University of Toronto Press.

\bibitem[\citeproctext]{ref-Fulk2018}
Fulk, R. D. 2018. \emph{A Comparative Grammar of the Early {Germanic}
Languages}. Vol. 3. Studies in Germanic Linguistics. John Benjamins.

\bibitem[\citeproctext]{ref-Hogg1992}
Hogg, Richard M. 1992. \emph{A Grammar of Old English. Volume 1:
Phonology}. Blackwell.

\bibitem[\citeproctext]{ref-Kroonen2013}
Kroonen, Guus. 2013. \emph{Etymological Dictionary of {Proto-Germanic}}.
Vol. 11. Leiden Indo-European Etymological Dictionary Series. Brill.

\bibitem[\citeproctext]{ref-Orel2003}
Orel, Vladimir. 2003. \emph{A Handbook of {Germanic} Etymology}. Brill.

\bibitem[\citeproctext]{ref-RingeTaylor2014}
Ringe, Don, and Ann Taylor. 2014. \emph{The Development of {Old
English}}. Vol. 2. A Linguistic History of English. Oxford University
Press.

\bibitem[\citeproctext]{ref-SieversBrunner1965}
Sievers, Eduard. 1965. \emph{Altenglische Grammatik}. 3rd edn. Edited by
Karl Brunner. Niemeyer.

\end{CSLReferences}

\end{document}
